% Source-aligned English for printed p.178.
\SourcePage{178}

\LayerTarget{tohoku:criterion:cech-H2-example-reduction:p2}

\LayerTarget{tohoku:example:cech-H2-counterexample-construction}
Let $X$ be an irreducible space (see the end of 3.3), and let $Y_1$ and $Y_2$
be two irreducible closed subsets of $X$ meeting in exactly two points $x_1$
and $x_2$ (for example, two intersecting circles in the plane endowed with the
Zariski topology); let $Y$ be their union. By abuse of notation, write
$\mathbf Z$ for the constant sheaf of integers on $X$, and consider, with the
notation of 3.5, the exact sequence of sheaves

\LayerTarget{tohoku:eq:closed-complement-constant-sheaf}
\[
0\longrightarrow\mathbf Z_{C_Y}\longrightarrow\mathbf Z
\longrightarrow\mathbf Z_Y\longrightarrow 0.
\]

\LayerTarget{tohoku:claim:cech-H2-counterexample-sheaf}
I claim that the sheaf $F=\mathbf Z_{C_Y}$ satisfies the required conditions.
First, since $H^p(X,\mathbf Z)=0$ for $p>0$ by the end of 3.3, we have
$H^2(X,\mathbf Z_{C_Y})=H^1(X,\mathbf Z_Y)=H^1(Y,\mathbf Z)$. Let us show
that this group is nonzero and, more precisely, isomorphic to $\mathbf Z$.
Indeed, there is a natural monomorphism from the constant sheaf $\mathbf Z$ on
$Y$ into the direct sum of the sheaves $\mathbf Z_{Y_1}$ and
$\mathbf Z_{Y_2}$, giving an exact sequence of sheaves on $Y$:

\LayerTarget{tohoku:eq:union-components-constant-sheaves}
\[
0\longrightarrow\mathbf Z_Y
\longrightarrow(\mathbf Z_{Y_1}+\mathbf Z_{Y_2})
\longrightarrow\mathbf Z_{Y_1\cap Y_2}\longrightarrow 0.
\]

\LayerTarget{tohoku:calculation:H1-union-cokernel}
Since $Y_i$ is irreducible, we again have
$H^p(Y,\mathbf Z_{Y_i})=H^p(Y_i,\mathbf Z)=0$ for $p>0$, whence
\[
H^1(Y,\mathbf Z)
=\Gamma(\mathbf Z_{Y_1\cap Y_2})/
\operatorname{Im}\bigl(\Gamma(\mathbf Z_{Y_1}+\mathbf Z_{Y_2})\bigr).
\]
This is the cokernel of the homomorphism
$\mathbf Z^2\longrightarrow\mathbf Z^2$ given by
$(n_1,n_2)\longmapsto(n_1-n_2,n_1-n_2)$, hence is isomorphic to $\mathbf Z$.
Thus

\LayerTarget{tohoku:eq:counterexample-H2-is-Z}
\[
H^2(X,\mathbf Z_{C_Y})=H^1(X,\mathbf Z_Y)=\mathbf Z.
\]

\LayerTarget{tohoku:goal:counterexample-cech-H1-is-Z}
It remains to prove that $\check H^1(X,H^1(\mathbf Z_{C_Y}))$ is isomorphic to
$\mathbf Z$ (for then, in view of the preceding relation, the epimorphism
$H^2(X,F)\longrightarrow\check H^1(X,H^1(F))$ must be an isomorphism). Let us
compute $H^1(\mathbf Z_{C_Y})$. For every open set $V$,
$H^1(V,\mathbf Z_{C_Y})$ can be computed from the following exact sequence of
sheaves on $V$:

\LayerTarget{tohoku:eq:restricted-complement-constant-sheaf}
\[
0\longrightarrow\mathbf Z_{C_Y}\longrightarrow\mathbf Z
\longrightarrow\mathbf Z_{Y'}\longrightarrow 0,
\]
where $Y'=Y\cap V$. Since $V$ is also irreducible, it follows that

\LayerTarget{tohoku:eq:local-H1-reduced-H0}
\[
H^1(V,\mathbf Z_{C_Y})
=H^0(Y',\mathbf Z)/\operatorname{Im}H^0(V,\mathbf Z)
=\widetilde H^0(Y',\mathbf Z),
\]
where the last group denotes reduced integral cohomology in degree 0: here it
is the free abelian group generated by the connected components of $Y'$,
modulo the diagonal subgroup.

\LayerTarget{tohoku:criterion:local-H1-two-components}
Now $Y'=Y\cap V$ is an open subspace of $Y=Y_1\cup Y_2$, and therefore has
0, 1, or 2 connected components. The last case occurs exactly when $V$ meets
both $Y_1$ and $Y_2$ without meeting their intersection. Thus
$H^1(V,\mathbf Z_{C_Y})=0$ except in this last case.

\LayerTarget{tohoku:construction:special-point-indexed-coverings}
To compute
\[
\check H^1(X,H^1(F))=\varinjlim H^1(\mathcal U,H^1(F)),
\]
we may restrict to coverings $\mathcal U=(U_x)_{x\in X}$ such that each $U_x$
meets at most one of the two closed sets $Y_1$ and $Y_2$, except when $x$ is
one of the two points $x_1,x_2$ of $Y_1\cap Y_2$, in which case we require
$U_x$ not to contain the other point.

\LayerTarget{tohoku:eq:special-cover-C0-vanishing}
For such a $\mathcal U$ we immediately have $C^0(\mathcal U,H^1(F))=0$, so
$H^1(\mathcal U,H^1(F))$ identifies with the group
$Z^1(\mathcal U,H^1(F))$ of 1-cocycles of $\mathcal U$ with coefficients in
$H^1(F)$, say $(f_{x,y})_{x,y\in X}$.

\LayerTarget{tohoku:eq:special-cover-intersection-sections}
But
\[
\Gamma(U_x\cap U_y,H^1(F))=0
\]
unless $x=x_1$, $y=x_2$, or $x=x_2$, $y=x_1$. Hence

\LayerTarget{tohoku:eq:special-cover-C1-is-Z2}
\[
C^1(\mathcal U,H^1(F))\simeq\mathbf Z^2,
\]
and the cocycles identify with the pairs $(n,-n)$ of opposite integers, where
$n=f_{x_1,x_2}$.

\LayerTarget{tohoku:eq:special-cover-H1-is-Z}
Thus $H^1(\mathcal U,H^1(F))=\mathbf Z$, and hence, on passing to the limit,
$\check H^1(X,H^1(F))=\mathbf Z$, completing the proof.

\TohokuContinuation{179}
